\section*{3 Pre-processing Steps and Algebraic Tool Kit}
Before we start the Hensel lifting process, a polynomial should fulfill certain properties that the input polynomial might not have. In this section, we describe transformations of a polynomial that achieve these properties such that ABPs can compute the transformation and its inverse, and factors of the polynomials are maintained.
We also explain how to compute homogeneous components and how to solve linear systems via ABPs. We show how handle the special case when the given polynomial is just a power of an irreducible polynomial.
\subsection*{3.1 Computing homogeneous components and coefficients of a polynomial}
Let $p(x,\mathbf z)$ be polynomial of degree $d$ in variables $x$ and $\mathbf z=(z_1,\ldots,z_n)$. Let $B_p$ be an ABP of size $s$ that computes a polynomial $p$. Write $p$ as a polynomial in $x$, with coefficients from $\mathbb F[\mathbf z]$,
\[p(x,\mathbf z)=\sum_{i=0}^dp_i(\mathbf z)x^i.\]
We show that all the coefficients $p_i(\mathbf z)$ have ABPs of size $\operatorname{poly}(s,d)$.
The argument is similar to Strassen's homogenization technique for arithmetic circuits, an efficient way to compute all the homogeneous components of a polynomial. The same technique can be used for ABPs (see~\cite[Lemma 5.2 and Remark]{ref17}). Here we sketch the proof idea.
Each node $v$ of $B_p$ we split into $d+1$ nodes $v_0,\ldots,v_d$, such that node $v_i$ computes the degree $i$ part of the polynomial computed by node $v$, for $i=0,1,\ldots,d$. Consider an edge $e$ between node $u$ and $v$ in $B_p$.
\begin{itemize}\item If $e$ is labeled with a constant $c\in\mathbb F$ or a variable $z_i$, then we put an edge between $u_i$ and $v_i$ with label $c$ or $z_i$, respectively.\item If $e$ is labeled with variable $x$, then we put an edge between $u_i$ and $v_{i+1}$ with label1.\end{itemize}
The resulting ABP has one source node and $d+1$ sink nodes. The $i$-th sink node computes $p_i(\mathbf z)$.
For each edge of $B_p$ we get either $d$ or $d+1$ edges in the new ABP. Hence, its size is bounded by $s(d+1)$.
\setcounter{nttheorem}{3}
\begin{ntlemma}[Author Lemma 4: Coefficient extraction]
Let $p(x,\mathbf z)=\sum_{i=0}^dp_i(\mathbf z)x^i$ be a polynomial. Then $\sz(p_i)\leq(d+1)\sz(p)$, for $i=0,1,\ldots,d$.
\end{ntlemma}
The technique can easily be extended to constantly many variables. For two variables, consider $p(x,y,\mathbf z)=\sum_{i,j}p_{i,j}(\mathbf z)x^iy^j$. Then, from an ABP of size $s$ for $p$ we get ABPs for the coefficients $p_{i,j}(\mathbf z)$ of size $s(d+1)^2$ similarly as above.
In homogenization, we want to compute the homogeneous components of $p$. That is, write $p(\mathbf z)=\sum_{i=0}^dp_i(\mathbf z)$, where $\deg(p_i)=i$. From an ABP $B_p$ for $p$ we get ABPs for the $p_i$'s similarly as above: In the definition of the new edges, only for constant label, we put the edge from $u_i$ to $v_i$. In case of any variable label $z_j$, we put the edge from $u_i$ to $v_{i+1}$ with label $z_j$. Then the $i$-th sink node computes $p_i(\mathbf z)$. The size is bounded by $s(d+1)$.
\begin{ntlemma}[Author Lemma 5: Homogenization]
Let $p(\mathbf z)=\sum_{i=0}^dp_i(\mathbf z)$ be a polynomial with $\deg(p_i)=i$, for $i=0,1,\ldots,d$. Then $\sz(p_i)\leq(d+1)\sz(p)$, for $i=0,1,\ldots,d$.
\end{ntlemma}
\subsection*{3.2 Computing $q$ from $p=q^e$}
A special case is when the given polynomial $p(\mathbf x)$ is a power of one irreducible polynomial $q(\mathbf x)$, i.e.,$p=q^e$, for some $e>1$. This case is handled separately. Kaltofen~\cite{ref10} showed how to compute $q$ for circuits, ABPs, and arithmetic formulas. Here, we give a short proof from Dutta~\cite{ref5}.
\begin{ntlemma}[Author Lemma 6]
Let $p=q^e$, for polynomials $p(\mathbf x),q(\mathbf x)$. Then $\sz(q)\leq\operatorname{poly}(\sz(p))$.
\end{ntlemma}
\begin{proof}
We may assume that $p$ is nonzero; otherwise the claim is trivial. We want to apply Newton's binomial theorem to compute $q=p^{1/e}$. For this we need that $p(0,\ldots,0)=1$. If this is not the case, we first transform $p$ as follows.
\begin{enumerate}\item If $p(0,\ldots,0)=0$, let $\boldsymbol\alpha=(\alpha_1,\ldots,\alpha_n)$ be a point where $p(\boldsymbol\alpha)\ne0$. By the PIT-Theorem, a random point $\boldsymbol\alpha$ will work, with high probability. Now we shift the variables and work with the shifted polynomial $\widetilde p(\mathbf x)=p(\mathbf x+\boldsymbol\alpha)$. Still,$\widetilde p(0,\ldots,0)$ might be different from1. In this case, we also apply the next item to $\widetilde f$.\item If $p(0,\ldots,0)=a_0\ne0$, then we work with $\widetilde p(\mathbf x)=p(\mathbf x)/a_0$. Then $\widetilde p(0,\ldots,0)=1$.\end{enumerate}
Note that both transformations are easily reversible. Hence, in the following we simply assume that $p(0,\ldots,0)=1$.
By Newton's binomial theorem, we have
\begin{equation}q=p^{1/e}=(1+(p-1))^{1/e}=\sum_{i=0}^\infty\binom{1/e}i(p-1)^i.\tag{1}\end{equation}
Note that $p^{1/e}$ is a polynomial of degree $d=\deg(q)$. Since $p-1$ is constant free, the terms $(p-1)^j$ in the RHS of(1) have degree $>d$, for $j>d$. Thus(1) turns into a finite sum modulo the ideal $\langle\mathbf x\rangle^{d+1}$,
\begin{equation}q=\sum_{i=0}^d\binom{1/e}i(p-1)^i\pmod{\langle\mathbf x\rangle^{d+1}}.\tag{2}\end{equation}
Let $\sz(p)=s$. For the polynomial $Q=\sum_{i=0}^d\binom{1/e}i(p-1)^i$ from(2), we clearly have $\sz(Q)\leq\operatorname{poly}(s)$. Finally, to get $q=Q\bmod\langle\mathbf x\rangle^{d+1}$, we have to truncate the terms in $Q$ with degree $>d$. This can be done by computing the homogeneous components of $Q$ as described in Lemma 5. We conclude that $\sz(q)\leq\operatorname{poly}(s)$.
\end{proof}
\subsection*{3.3 Reducing the multiplicity of a factor}
In the earlier works on bivariate and multivariate polynomial factoring, typically the problem is reduced to factoring a square-free polynomial. This is convenient at various places in the Hensel lifting process. The technique to reduce to the square-free case is via taking the gcd of the input polynomial and its derivative. However, for getting upper bounds on the ABP-size of the factors, we want to avoid gcd-computations, because no polynomial size upper bound for the gcd of two ABPs is known.
We avoid this problem by observing that we do not need the polynomial to be square-free. As we will see, it suffices to have one irreducible factor with multiplicity one, and another coprime factor.
Let $p(\mathbf x)$ be the given polynomial, for $\mathbf x=(x_1,\ldots,x_n)$. The special case that $p$ is a power of one irreducible polynomial we just handled in Section 3.2. Hence, we may assume that $p$ has at least two irreducible factors. So let $p=q^ep'$, where $q$ is irreducible and coprime to $p'$.
Consider the derivative of $p$ w.r.t. some variable, say $x_1$.
\begin{equation}\frac{\partial p}{\partial x_1}=q^{e-1}\left((e-1)\frac{\partial q}{\partial x_1}p'+q\frac{\partial p'}{\partial x_1}\right).\tag{3}\end{equation}
Note that $q$ does not divide the factor $(e-1)\frac{\partial q}{\partial x_1}p'+q\frac{\partial p'}{\partial x_1}$ in(3). Hence, the multiplicity of factor $q$ in $\frac{\partial p}{\partial x_1}$ is reduced by one compared to $p$.
For the ABP-size, we write $p$ as a polynomial in $x_1$, i.e.$p(\mathbf x)=\sum_{i=0}^da_ix_1^i$, where the coefficients $a_i$ are polynomials in $x_2,\ldots,x_n$. By Lemma 4, when $p$ has an ABP of size $s$, then the coefficients $a_i$ can be computed by ABPs of size $s'=s(d+1)$. We observe that then the coefficients of the derivative polynomial $\frac{\partial p}{\partial x_1}=\sum_{i=1}^d ia_ix_1^{i-1}$ have ABPs of size $s'+1$.
We repeat taking derivatives $k=e-1$ times and get $\frac{\partial^kp}{\partial x_1^k}$, which has the irreducible factor $q$ with multiplicity one, as desired.
The coefficients of $\frac{\partial^kp}{\partial x_1^k}$ can be computed by ABPs of size $s'+1$. This yields an ABP of size $\operatorname{poly}(s)$ that computes $\frac{\partial^kp}{\partial x_1^k}$.
